\subsection{局部凸空间的定义}

定义 1. --- 拓扑向量空间称为局部凸的（实的），如果存在由凸集组成的 0 的一个基本邻域系。

这样的空间称为局部凸空间。它的拓扑称为局部凸拓扑。

我们在本书以下各部分所研究的 \(\mathbf{R}\) 上的拓扑向量空间几乎都是局部凸的。

若 V 是局部凸空间 E 中 0 的一个凸邻域，则 \(V \cap (-V)\) 是 0 的一个对称凸邻域。由于凸集的闭包是凸集 (II, 第 13 页, prop. 14)，由 I, 第 7 页, prop. 4 可知，E 中闭、对称且凸的那些 0 的邻域构成一个在以 0 为中心、比 \(\neq 0\) 的位似下不变的基本邻域系。

命题 1. --- 设 S 是向量空间 E 上由吸收、对称且凸的集组成的一个滤子基。则 S 中的集经比 \textgreater{} 0 的位似变换所得的集所成的 V，构成 E 上某个局部凸拓扑的 0 的一个基本邻域系。

显然 \(\mathfrak{B}\) 是满足 I, 第 7 页, prop. 4 中条件 \((\mathrm{EV}_{\mathrm{I}})\) 与 \((\mathrm{EV}_{\mathrm{II}})\) 的滤子基；它也满足 \((\mathrm{EV}_{\mathrm{III}})\) ，因为若 \(\mathbf{V} \in \mathfrak{S}\) ，则 \(\frac{1}{2}\mathbf{V} + \frac{1}{2}\mathbf{V} = \mathbf{V}\) 。

注意，若 \(\mathcal{T}\) 是 E 上以 \(\mathfrak{V}\) 为 0 的基本邻域系的局部凸拓扑，则诸集 \((1/n)\) V（其中 \(n\) 取遍整数 \(>0\) 且 V 取遍 \(\mathfrak{S}\) ）构成拓扑 \(\mathcal{T}\) 的 0 的一个基本邻域系。\(\mathcal{T}\) 是 Hausdorff 的，必须且只需对 E 中每个 \(x \neq 0\) ，存在一个整数 \(n\) 及一个集 \(\mathrm{V} \in \mathfrak{S}\) ，使 \(nx \notin \mathrm{V}\) ；此外，若 \(\mathfrak{S}\) 可数，则拓扑 \(\mathcal{T}\) 是可度量化的局部凸拓扑。反之，显然若 \(\mathcal{T}\) 是可度量化的局部凸拓扑，则存在 \(\mathcal{T}\) 的由闭、对称、凸的 0 邻域组成的一个可数基本邻域系。

推论. --- 拓扑向量空间 E 的拓扑 \(\mathcal{T}\) 由一族半范数定义 (II, 第 3 页)，必须且只需 \(\mathcal{T}\) 是局部凸的。

条件是必要的，因为 E 上每个半范数都是凸函数，从而对 \(\alpha > 0\) ，\(x \in E\) 中使 \(p(x) \leqslant \alpha\) 的点所成的集是凸集 (II, 第 17 页, 推论)。反之，若 V 是 E 中 0 的对称、闭、凸邻域，则 V 的度规 p 是 E 上的半范数，使 V 恰为 E 中满足 \(p(x) \leqslant 1\) 的点 x 所成的集 (II, 第 20 页, prop. 23)。

这进一步表明，局部凸拓扑 \(\mathcal{T}\) 由对 \(\mathcal{T}\) 连续的全体半范数所成的族定义。此外，若 \(\mathcal{T}\) 可度量化，则它由可数个半范数定义。

由 prop. 1 的推论，§ 1 中关于由半范数族定义的拓扑的全部结果，都特别适用于实向量空间上的局部凸拓扑。局部凸 Hausdorff 空间 E 有一个局部凸的完备化 \(\hat{E}\) 。完备可度量的局部凸空间称为 Fréchet 空间；每个 Banach 空间都是 Fréchet 空间。

命题 2. --- 设 f 是局部凸空间 E 的向量子空间 M 上定义的连续线性型；则存在 E 上定义的连续线性型 h，它是 f 的扩张。

由上面的推论及 II, 第 7 页, cor. 2，存在连续半范数 \(p\) 定义在 \(\mathbf{E}\) 上，使得 \(|f(y)| \leqslant p(y)\) 对所有 \(y \in \mathbf{M}\) 成立。由 Hahn-Banach 定理 (II, 第 23 页, cor. 1)，存在线性型 \(h\) 定义在 \(\mathbf{E}\) 上，它扩张 \(f\) 并且 \(|h(x)| \leqslant p(x)\) 对所有 \(x \in \mathbf{E}\) 成立，这就蕴含 \(h\) 连续 (II, 第 6 页, prop. 5)。

注记. --- 若 g 是 M 到乘积空间 \(R^{I}\) 中的连续线性映射，则存在 E 到 \(R^{I}\) 中的连续线性映射 h，它是 g 的扩张；因为把 g 写作 \(g = (g_{\mathrm{t}})\) ，其中诸 \(g_{t}\) 是 M 上定义的连续线性型，则存在 \(h_{t}\) ，它是 \(g_{t}\) 对每个 \(t \in I\) 的一个扩张，使 \(h_{t}\) 是 E 上的连续线性型。连续线性映射 \(h = (h_{\mathrm{t}})\) 即具有所要求的性质。

注意，若 \(\mathbf{F}\) 是局部凸 Hausdorff 空间，而 \(g\) 是 \(\mathbf{M}\) 到 \(\mathbf{F}\) 中的连续线性映射，则未必存在 \(\mathbf{E}\) 到 \(\mathbf{F}\) 中的、扩张 \(g(\mathrm{IV},\mathrm{p.~55},\mathrm{exerc.~16},c))\) 的连续线性映射。然而当 \(\mathbf{M}\) 有限维时，这样的扩张确实存在（见下文 cor. 2）。

推论 1. --- 设 E 是局部凸空间。若 \(x_0 \in E\) 不属于 \(\{0\}\) 的闭包，则存在 E 上定义的连续线性型 \(f\) ，使 \(f(x_0) \neq 0\) 。

对 \(x_0\) 生成的一维向量空间 M 以及 M 上定义的线性型 \(\xi x_0 \mapsto \xi\) 应用 prop. 2；由 I, 第 12 页, prop. 2，该线性型是连续的。

推论 2. --- 设 M 是局部凸 Hausdorff 空间 E 的有限维向量子空间。则存在 E 的闭向量子空间 N，它是 M 在 E 中的拓扑补。

M 在 E 中存在拓扑补，必须且只需 M 到自身的恒等映射可以扩张为 E 到 M 上的连续线性映射，而这个映射此时必然是一个连续投影算子 (GT, III, § 6.2, 推论)。而由上面的注记即可得到这一点，因为 M 同构于空间 \(R^{n}\) (I, 第 13 页, th. 2)。

命题 3. --- 在局部凸空间 E 中，预紧集的均衡凸包本身是一个预紧集。

设 A 是 E 中的预紧集。给定 E 中 0 的均衡凸邻域 V，存在有限多个点 \(a_i \in \mathbf{A}\) （ \(1 \leqslant i \leqslant n\) ），使 A 含于诸邻域 \(a_i + \mathbf{V}\) （ \(1 \leqslant i \leqslant n\) ）的并 S 中。于是 A 的均衡凸包 C 含于 S 的均衡凸包 T；而 T 含于 B + V，其中 B 表示有限点集 \(a_i\) ， \(-a_i\) （ \(1 \leqslant i \leqslant n\) ）的凸包。又 B 是预紧的 (II, 第 14 页, cor. 2)；故存在有限多个点 \(b_k \in \mathbf{B}\) （ \(1 \leqslant k \leqslant m\) ），使 \(\mathbf{B}_k\) 含于诸邻域 \(b_k + \mathbf{V}\) 的并中。于是 C 含于诸邻域 \(b_k + 2\mathbf{V}\) 的并中，命题得证。

注意，在无限维局部凸 Hausdorff 空间中，紧集的凸包未必是闭的 (II, 第 74 页, exerc. 3)。

推论. --- 若在局部凸 Hausdorff 空间 E 中，紧集 X 含于一个完备凸集中（完备是就对 E 诱导的一致结构而言），则 X 的闭凸包是紧的。

因为这个包是完备空间的一个闭子集，因而是完备的，同时它又是预紧的且 Hausdorff 的。

然而在非完备的局部凸 Hausdorff 空间中，紧集的闭凸包未必是紧的 (II, 第 87 页, exerc. 2)。

